% Candidate replacing the safe-common-edge proof and its interface table.
% Existing cycles and safe-edge definitions remain unchanged.
% No main-manuscript edit. All colour symbols below are distinct.
\paragraph{Cycle deletion and common edges.}
For cycles $C,D$, let $U=V(C)\cap V(D)\ne\varnothing$ and let $H$
contain their edges with endpoints in $U$. Let $F\subseteq E(H)$ be
forbidden, and let $R\subseteq E(H)\setminus F$ be used by neither
cycle. If deleting the outside vertices removes $q_C,q_D$ paths,
counting isolated vertices as paths, then
\[
 |(E(C)\cap E(D))\setminus F|
 \ge 2|U|-q_C-q_D-|E(H)|-|F|+|R|.                 \tag{DP}
\]
Indeed, either cycle has $|U|-q$ edges with endpoints in $U$, with
$q=0$ for an empty outside set. Deleting forbidden edges loses at
most $|F|$ from each. Both resulting edge sets lie in
$E(H)\setminus(F\cup R)$, so inclusion--exclusion proves (DP).

\paragraph{A proof of the safe-common-edge lemma.}
Write the good boundary types as $E(aab),M(aba),A(abc)$, allowing
reversal. Orient an interface from the smaller type in $E<M<A$.
From $aab$, changing the singleton gives $aac$; changing the middle
gives $abb$ or $acb$; changing the repeated endpoint gives $bab$ or
$cab$. From $aba$, the remaining changes give $aca$ or $cba$.
From $abc$, they give $abd$ or $adc$, after reversal if needed.
Thus there are nine families, where $a,b,c,d$ are distinct.

For four families use (DP), with $H$ the common recolouring graph and
$F=(P_t\cup P_s)\cap E(H)$:
\[
\begin{array}{c|rrrrr|r}
 t\to s&|U|&|E(H)|&|F|&q_C+q_D&|R|&\text{bound}\\\hline
 aab\to aac&10&13&2&2&0&3\\
 aab\to acb&10&15&1&2&0&2\\
 aab\to cab&9&11&1&4&0&2\\
 abc\to abd&11&17&0&4&2&3
\end{array}                                                     \tag{CT}
\]
The fixed-middle layers in these common graphs are respectively
$(C_4,C_4,K_2)$, $(K_2\square K_3,C_4)$, $(K_3,C_4,K_2)$, and
$(K_3,C_4,C_4)$, with four, two, three, and six middle-change edges.
These give the vertex and edge counts. The two special edges of $E(aab)$
are $bcd$--$dcd$ and $bdc$--$cdc$: one from each $E$ side survives in
the first row, and one from the source survives in either following row.
The deletion identities below give the values of
$q_C+q_D$. For the last row, both cycles omit $bab$--$cab$ and
$bab$--$dab$ in the middle-$a$ triangle: their middle-$a$ paths start
at $bab$ by changing the last coordinate to the respective unused
colour. These two omitted edges form $R$.

For $aab\to abb$, the common middle-$c$ square has opposite corners
$bca,dcd$. Both cycles retain its two edges at $bca$. Their special
edges lie at $dcd$ and are absent from the other cycle, so the two
common edges are safe. For $abc\to adc$, both cycles retain five
edges of the common middle-$a$ hexagon
\[
(bab,bad,cad,dad,dab,cab,bab);
\]
hence they share at least four. For $aba\to aca$, the preceding
square-support property gives three retained edges out of the four
in the middle-$d$ square $Q_d\subset B_{aaa}$, hence two common edges.
The latter two families have no special edges.

Only $aab\to bab$ and $aba\to cba$ remain. Let $M_0(aba)$ be the
image of $D_M$ under $(0,1,2,3)\mapsto(a,b,c,d)$, and $M_1(aba)$
its reversal. Unused-colour interchange fixes both edge sets.
The following four pairs follow directly from \eqref{eq:path232-cycles}:
\[
\begin{array}{c|c|l}
\text{family}&\text{M orientation}&\text{two safe common edges}\\\hline
aab\to bab&0&cba\text{--}cbd,\ dba\text{--}dbc\\
aab\to bab&1&cbc\text{--}dbc,\ cbd\text{--}dbd\\
aba\to cba&0&dab\text{--}dac,\ dac\text{--}dad\\
aba\to cba&1&bcb\text{--}bcd,\ dcb\text{--}dcd
\end{array}                                                     \tag{MP}
\]
The orientation belongs to the target in the first family and to the
source in the second. In the first family neither $E$ special edge
belongs to the common graph; the second has no special edges.
The $E$ and $A$ edge sets at a fixed boundary are unique: unused-colour
interchange fixes $E$, and reversal followed by endpoint-colour
interchange fixes $A$. Thus the proof covers every allowed cycle.
Colour permutations and reversal preserve the special-edge sets and
transport every argument to all labelled interfaces.

\paragraph{Six deletion identities supporting (CT).}
Use $[x_0,\ldots,x_k]$ for a path with the stated consecutive vertices.
In the $E(aab)$ cycle, the classes $w=c$, $v=c$, and $u=c$ respectively
form
\[
[dbc,cbc,cdc,bdc],\quad[dca,bca,bcd,dcd],\quad
[cbd,cba,cda]\ \cup\ [cbc,cdc].
\]
In the $A(abc)$ cycle, the classes $v=a$, $u=b$, and $w=d$ respectively
form
\[
[bab,bad,cad,dad,dab,cab],\quad
[bad,bab,bdb,bcb]\ \cup\ [bcd,bca,bda],\quad
[bad,cad,dad]\ \cup\ [dcd,bcd].
\]
Legality lists give all the vertices in each class, and
\eqref{eq:path232-cycles} gives exactly the displayed path edges.
Interchanging unused colours preserves $E$; the reversal symmetry
preserves $A$. Thus these identities apply to every allowed cycle,
independently of its colour identification. Substitution gives the
two path counts in each row of (CT).

% Mathematical boundary: DP is general, but the deletion identities,
% square/hexagon supports and MP are properties of the selected cycles.
% No arbitrary-Hamilton-cycle claim is made for path232.
% This candidate removes the target-convention/orbit bookkeeping and
% six of ten literal safe-edge entries; four parametric entries remain.
